% appendix.tex
% Extended prose notes on prior work. The main Existing Research section
% (existing_research.tex) summarizes these systems in tables; this appendix
% preserves the per-system detail for readers who want it.

\appendix

\section{Extended Notes on Prior Work}
\label{sec:appendix_prior_work}

This appendix expands the tabular summaries in
Section~\ref{sec:existing_research} into per-system prose. It adds no new
claims; it records the same systems in more detail for readers who want it.

\subsection{Research Infrastructure}
\label{subsec:app_infrastructure}

Almost all modern Pok\'{e}mon AI research is built on \emph{Pok\'{e}mon
Showdown}, an open-source, community-maintained battle simulator that
reimplements the game's mechanics and runs online ladders for nearly every
format~\cite{pokemon_showdown}. Showdown matters to research for practical
reasons: it executes battles far faster than cartridge play, it exposes machine-
readable battle logs and replays, and it supports a wide range of historical and
current formats, including doubles. It is an approximation of the official games
rather than a certified copy of them, but it is the de facto environment for the
field.

The standard programmatic interface to Showdown is \emph{poke-env}, a Python
library that lets an agent connect to a Showdown server, read the battle state,
and submit actions~\cite{poke_env}. It underpins most of the systems discussed
below and is the usual starting point for reinforcement-learning experiments,
scripted baselines, and self-play data collection. The library was developed
primarily for singles play; support for doubles and VGC-specific structure has
been added more recently and by individual projects rather than as a mature,
standardized capability.

The use of Showdown as an AI benchmark was established early by the
\emph{Showdown AI Competition} of Lee and Togelius~\cite{lee2017showdown}, which
framed competitive Pok\'{e}mon singles as a research challenge and built
infrastructure for automated agents to play it. More recently, the
\emph{PokeAgent Challenge}~\cite{karten2025pokeagent}, run as a NeurIPS~2025
competition, has positioned Pok\'{e}mon as a large-scale benchmark for
decision-making under partial observability, game-theoretic reasoning, and
long-horizon planning. It comprises two tracks, a battling track on Showdown and
a single-player speedrunning track, supplies a dataset of over 20 million battle
trajectories and a suite of baseline agents, and drew over one hundred competing
teams. Its battling track is centered on singles formats, which is itself a useful
signal: even the field's flagship community benchmark treats singles as the
primary battling target.

\subsection{Competitive Pok\'{e}mon Singles: State of the Art}
\label{subsec:app_singles}

The strongest published battle results in competitive Pok\'{e}mon are in the
singles format. Three systems define the current state of the art, and together
they map out the two methodological families (large-scale learning from human
data, and search guided by learned or language-model priors) that a VGC study is
most likely to inherit.

\paragraph{Offline reinforcement learning from human replays (Metamon).}
The system closest in spirit to a deep-learning approach is Metamon, presented by
Grigsby et al. in \emph{Human-Level Competitive Pok\'{e}mon via Scalable Offline
Reinforcement Learning with Transformers}~\cite{grigsby2025metamon}. Its central
idea is to reconstruct an agent's first-person trajectory from third-person
spectator replays accumulated over more than a decade of human play, then train
Transformer policies of up to 200M parameters on that data using imitation
learning and offline reinforcement learning. Evaluated in the four oldest and
most partially observed Pok\'{e}mon generations (Gen~1--4), the best agents climb
into the top 10\% of active human players on the Showdown ladder. Methodologically,
this is the most direct existing template for learning a Pok\'{e}mon battle policy
at scale from historical human data rather than from self-play alone. The work
targets singles formats.

\paragraph{Search with language-model priors (Pok\'{e}Champ).}
Pok\'{e}Champ, by Karten et al., takes a different route~\cite{karten2025pokechamp}.
It runs an explicit minimax tree search over battle states, but replaces the
three components that such a search normally needs (action proposal, opponent
modeling, and value estimation) with a large language model, so that the agent
requires no Pok\'{e}mon-specific training. It reports strong win rates against
prior language-model and rule-based agents in the Gen~9 OU singles format and, in
the process, compiles a large dataset of real human games. Pok\'{e}Champ is
evidence for a hybrid architecture (search combined with learned or model-based
priors and opponent modeling) as a credible alternative to end-to-end
reinforcement learning. It is evaluated in singles.

\paragraph{Language-model agents (Pok\'{e}LLMon).}
Pok\'{e}LLMon, by Hu et al., was the first language-model agent to reach human
parity in Pok\'{e}mon battles~\cite{hu2024pokellmon}. It combines in-context
learning, knowledge-augmented generation, and mechanisms to keep its chosen
actions consistent, and reaches roughly an even win rate on the Showdown ladder.
It is foundational for the language-model line of work that Pok\'{e}Champ later
built on, though it is more modest in strength. It targets singles.

\paragraph{Earlier foundations.}
Before this recent wave, self-play reinforcement learning and search-based agents
were applied to Pok\'{e}mon singles on Showdown, growing out of the benchmark that
Lee and Togelius established~\cite{lee2017showdown}. These earlier systems are
mainly of value as baselines and as evidence that Pok\'{e}mon battling has been an
active AI target for some years; they did not reach the strength of the systems
above, and they too are singles work.

\subsection{Dedicated VGC (Doubles) Work}
\label{subsec:app_vgc}

A smaller body of work targets the VGC doubles format directly. It is recent,
dedicated, and thinner than the singles literature, but it includes the most
important reference point for this study.

\paragraph{A benchmark for the real game (VGC-Bench).}
The anchor work is VGC-Bench, by Angliss et al., accepted to
AAMAS~2026~\cite{angliss2026vgcbench}. It is the first rigorous benchmark built
specifically for VGC doubles. It is implemented on top of poke-env (with a
PettingZoo multi-agent integration), ships a human-play dataset of more than
700{,}000 open-team-sheet VGC battle logs parsed from Showdown, and provides a
suite of baseline agents spanning heuristics, a
language-model agent, behavior cloning, and multi-agent reinforcement learning
trained under several empirical game-theoretic schemes (self-play, fictitious
play, and double oracle), including variants initialized from behavior cloning.
Its central empirical finding is a sharp contrast: an agent trained and evaluated
on a single fixed team can reach expert strength, to the point of defeating a
professional VGC competitor, yet performance degrades substantially as the set of
training teams grows, and even strong agents remain exploitable. VGC-Bench
studies team \emph{usage}: the construction of teams is left outside its scope.
This benchmark is the strongest single piece of evidence about where dedicated
VGC work currently stands, and it is cited heavily in the analysis that follows.

\paragraph{Abstracted VGC and metagame balance (the Reis--Lau line).}
The longest-running dedicated VGC research program comes from the Portuguese
groups at Porto and Aveiro. Their \emph{VGC AI Competition} defines a framework
with modules for team selection, battling, team building, and metagame balance,
and poses the production of a balanced metagame (one in which many configurations
remain viable) as an AI challenge~\cite{reis2021vgccompetition}. A later paper
extends this into an adversarial loop in which team-building agents push toward
the strongest configurations, collapsing diversity, while a balancing agent
re-tunes Pok\'{e}mon attributes to restore it~\cite{reis2023adversarial}. This
line is important for its explicit treatment of team building and the metagame as
objects of study. Its main caveat is that it operates over an abstracted model of
Pok\'{e}mon rather than the full game engine, so it speaks to the metagame and
team-building framing more than to battle performance in the real game.

\paragraph{Deep RL and neuroevolution on the real format.}
Rodriguez, Villanueva, and Bald\'{e}on apply deep reinforcement learning together
with neuroevolution to real VGC double battles, and also build a move-suggestion
tool and promote the use of underrepresented Pok\'{e}mon for
diversity~\cite{rodriguez2024neuroevolution}. It is one of the few works that
targets the genuine doubles format with learning methods, at a smaller scale than
VGC-Bench, and it confirms that the on-format problem has begun to attract direct
attention.

\subsection{Analogous Games and Their AI Systems}
\label{subsec:app_analogous}

VGC's difficulty comes from combining four strategic properties at once
(long-horizon planning, simultaneous action, imperfect information, and
stochasticity), as set out in Section~\ref{sec:introduction}. No single solved
game combines all four, but several landmark systems each solved a game that
shares one or more of them, and these are the systems whose methods a VGC study
would draw on. They are reviewed here as methodological precedent; the argument
that none of them covers VGC's full combination is developed in
Section~\ref{sec:analysis}.

\paragraph{Perfect-information self-play (AlphaGo and AlphaZero).}
AlphaGo first reached superhuman strength in Go by combining deep neural networks
with Monte Carlo tree search~\cite{silver2016alphago}, and AlphaZero generalized
the recipe to chess, shogi, and Go from self-play alone, with no human data and
no domain knowledge beyond the rules~\cite{silver2018alphazero}. This is the
canonical demonstration that self-play plus search can produce superhuman play,
and it is the natural first analogy for any new game. Its setting, however, is the
opposite of VGC's on three of the four pillars: Go and chess are
perfect-information, deterministic, and strictly turn-based. The transferable
lesson is the self-play-with-search paradigm itself, not the specific algorithm.

\paragraph{Imperfect information and search (poker: Libratus, NFSP, ReBeL).}
Poker is the canonical imperfect-information benchmark and is in several respects a
closer analogy than Go. Libratus reached superhuman play in heads-up no-limit
Texas hold'em, demonstrating that an equilibrium-seeking agent can beat top
professionals in a game of hidden information and chance~\cite{brown2017libratus}.
Neural Fictitious Self-Play (NFSP) showed that deep reinforcement learning can
approximate equilibrium strategies in imperfect-information games without
hand-crafted abstractions~\cite{heinrich2016nfsp}, and ReBeL unified
reinforcement learning with search for imperfect-information games, again reaching
superhuman heads-up poker~\cite{brown2020rebel}. The lesson this line carries for
VGC is that the right formal object is not a standard Markov decision process but
an imperfect-information game in which the goal is an equilibrium (an
unexploitable strategy), and that combining learning with search is a productive
way to reach it. Poker is, however, strictly turn-based, so it does not exercise
VGC's simultaneous-action structure.

\paragraph{Large imperfect-information game trees (Stratego: DeepNash).}
DeepNash reached expert human level in Stratego using model-free multi-agent
reinforcement learning driven by a regularized form of Nash dynamics, without any
explicit search~\cite{perolat2022deepnash}. Stratego is the closest cousin to VGC
among solved games on the information axis: it has a very large
imperfect-information game tree and rewards long-horizon deception. It shows that
robust, near-equilibrium policies can be learned directly, even when the game is
too large to search explicitly. It differs from VGC in being deterministic and
strictly alternating, so it isolates the imperfect-information pillar without the
stochastic and simultaneous-action pillars.

\paragraph{Population-based training at scale (StarCraft II: AlphaStar).}
AlphaStar reached grandmaster level in StarCraft II, a real-time game with a
very large action space and partial observability through fog of
war~\cite{vinyals2019alphastar}. Its most relevant contribution for VGC is
structural rather than architectural: the \emph{league}, a population of agents,
including dedicated ``exploiter'' agents that seek out and punish weaknesses in
the main agents, trained against one another. This population structure is a
response to the fact that naive self-play can overfit to its own strategies and
cycle in games where strategies are non-transitive, a property VGC exhibits
strongly. StarCraft II has long horizons and partial information but is essentially
deterministic and is dominated by real-time mechanical execution, the axis on
which VGC, being turn-based, places no demand.

\paragraph{Scaled self-play in a team game (Dota 2: OpenAI Five).}
OpenAI Five reached professional level in Dota 2, a long-horizon, partially
observable team game, through self-play reinforcement learning at very large
scale~\cite{berner2019openaifive}. It demonstrates that scaled self-play can work
in a complex competitive environment, while also illustrating the substantial
compute and engineering cost of doing so, a consideration any feasibility
assessment must weigh.

\paragraph{Population methods and simultaneous-move search.}
Two further strands of method are directly relevant to VGC's structure even though
they are not tied to a single flagship system. Policy-Space Response Oracles
(PSRO) and the broader empirical-game-theoretic and double-oracle family
generalize population training into a principled framework for games with
non-transitive strategy spaces~\cite{lanctot2017psro}; VGC-Bench's use of double
oracle and fictitious play sits within this family. Separately, work on Monte
Carlo tree search for simultaneous-move games adapts search to settings where both
players commit before either observes the other, using equilibrium approximation
rather than the ordinary one-player assumption that the opponent moves
second~\cite{lanctot2014smmcts}. Both strands speak directly to features of VGC
(its non-transitive metagame and its per-turn simultaneous commitment) that the
single-game flagship systems above do not individually address.

\newpage
